\documentclass[10pt]{ltjsarticle}
\usepackage{amsmath}
\usepackage{tikz}
\usetikzlibrary{matrix, positioning}
\usepackage{amssymb}
\usepackage{graphicx}
\usepackage{luatexja}
\usepackage{enumitem}
\setlist{itemsep=0pt, topsep=2pt, parsep=0pt}
\usepackage{luatexja-fontspec}
\usepackage{cite}
\usepackage{hyperref}
\usepackage[a4paper,top=10mm,bottom=12mm,left=18mm,right=18mm]{geometry}
\usepackage{listings}
\usepackage{xcolor}

%%% ========================================
\date{2026年08月03日}
\title{第14回リスクマネジメント課題}
\newcommand{\authorname}{学籍番号: 26720493\\山口紘太郎}
%%% ========================================
\setmainjfont[
  AutoFakeBold = false
]{IPAexMincho}

\setlength{\parskip}{2pt}
\linespread{0.95}

\makeatletter
\renewcommand{\maketitle}{%
  \begin{flushright}
    \small\@date
  \end{flushright}
  \vspace{0.1cm}
  \begin{center}
    {\large \@title}
  \end{center}
  \begin{flushright}
    {\normalsize \authorname}
  \end{flushright}
}
\makeatother

\begin{document}
{\small
\maketitle

\section{事故概要と原因}
本レポートで扱うのは，2020年（令和2年）12月4日に発生した，日本航空904便
（ボーイング777-200型，機体記号JA8978）の左側エンジン破損に関する事故である．
令和2年12月4日，那覇空港離陸後の上昇中，機体に振動を伴う異音が発生し，
機長はエンジンを停止させ那覇空港へ引き返した．着陸後の点検で，
ファンブレードの破断，カウリング等の一部脱落，胴体・水平尾翼の損傷が
確認されたが，乗員・乗客189名に負傷者はなかった．

原因は，ファンブレード製造時の研磨工程で溶着した微小な金属塊（ノジュール，
幅約0.12mm）を起点に亀裂が発生し，これが定期検査でも発見されないまま
運航が継続されたことで疲労破壊に至ったためと推定される．検査で亀裂が
発見されなかったのは，用いられた非破壊検査（TAI検査：Thermal Acoustic Image）の
手法・間隔が，亀裂発生部位の欠陥検出には不十分だったことによる．

\section{技術的側面の対策}
\subsection{自分の考え}
TAI検査は行われていたが，2回の検査でも亀裂は発見されなかった．微小な亀裂でも
重大事故につながりうるため，TAI検査の間隔（当時6500FCごと）を狭めることが
効果的な対策だと考える．ここで，FCとはフライト回数を意味する．

\subsection{実際の対策}
TAI検査の間隔を6500FCから1000FCへ
大幅短縮し，加えて275〜550FCごとにUT（超音波探傷）検査を導入した．
FAAも耐空性改善命令で，次の飛行までのTAI検査実施と，その後のTAI検査・UT
繰り返し実施を義務化した．

\section{経営上の対策}

\subsection{自分の考え}
亀裂の合否基準を明確にし，発見時は即使用停止とすることが重要である．
検査間隔の短縮も踏まえ，コストより安全を優先する姿勢が大事だと考える．

\subsection{実際の対策}
TAI検査の合否判断は個人ではなくチームで検討する方式に改訂された．FAAは
ほぼゼロ許容に近い厳格な運用へ移行し，国交省は対象機の運航停止を指示した．

\section{考察}
発見や事故発生がわずかでも遅れていれば，死傷者を伴う事故になっていた可能性が
あり，今回無傷であったのは結果論に過ぎないと考えるべきである．
実際の対策を見ると，技術的な検出限界を克服したのではなく，検査の頻度と判断基準を
引き上げたことで，「見逃しの確率を下げた」に近い．
TAI検査は非破壊検査として有効だが，検出が難しい部分が存在し，リスク管理にも限界が
あることが，この事故によって示されたといえる．

事故当時の時点でも，今回行われた対策は実現できたものも多い．検査期間の短縮や
合否判定なども即時対策ができた．
また，非破壊検査では，コストがかかるかもしれないが，TAI検査のみの検査ではなく，複数の検査技術を用いて
検査を行い，単一技術を過信しないことが重要であると感じた．
\begin{thebibliography}{9}

\bibitem{jtsb2022}
運輸安全委員会,
「航空重大インシデント調査報告書 AI2022-5 --- 日本航空株式会社所属 ボーイング式777-200型 JA8978 発動機の破損（破片が当該発動機のケースを貫通した場合に限る）に準ずる事態」,
2022年8月5日議決.
\url{https://jtsb.mlit.go.jp/aircraft/rep-inci/AI2022-5-1-JA8978.pdf}

\end{thebibliography}

\end{document}